\section{Constructible sets}
\label{section:0.9}

\subsection{Constructible sets}
\label{subsection:0.9.1}

\begin{definition}[9.1.1]
\label{0.9.1.1}
We say that a continuous map $f:X\to Y$ is \emph{quasi-compact} if
for every quasi-compact open subset $U$ of $Y$, $f^{-1}(U)$ is quasi-compact. We
say that a subset $Z$ of a topological space $X$ is \emph{retrocompact
in $X$} if the canonical injection $Z\to X$ is quasi-compact, in other words
if for every quasi-compact open subset $U$ of $X$, $U\cap Z$ is quasi-compact.
\end{definition}

A \emph{closed} subset of $X$ is retrocompact in $X$, but a
quasi-compact subset of $X$ is not necessarily retrocompact in $X$. If
$X$ is quasi-compact, every open subset retrocompact in $X$ is
quasi-compact. It is clear that every \emph{finite} union of sets
retrocompact in $X$ is retrocompact in $X$, every finite union
of quasi-compact sets being quasi-compact. Every finite intersection
of open subsets retrocompact in $X$ is an open subset retrocompact in $X$. In a
\emph{locally Noetherian} space $X$, every quasi-compact set is a
Noetherian subspace, and consequently \emph{every subset} of $X$ is
retrocompact in $X$.

\begin{definition}[9.1.2]
\label{0.9.1.2}
Given a topological space $X$, we say that a subset of $X$ is
\emph{constructible} if it belongs to the smallest set $\mathfrak{F}$ of
subsets of $X$ containing all the retrocompact open subsets of
$X$ and stable under finite intersection and passage to the complement
(which implies that $\mathfrak{F}$ is also stable under finite union).
\end{definition}

\begin{proposition}[9.1.3]
\label{0.9.1.3}
For a subset of $X$ to be constructible, it is necessary and sufficient
that it be a finite union of sets of the form $U\cap\complement{V}$, where $U$
and $V$ are open subsets retrocompact in $X$.
\end{proposition}

It is clear that the condition is sufficient. To see that it is
necessary, consider the set $\mathfrak{G}$ of finite unions
of sets of the form $U\cap\complement{V}$, where $U$ and $V$ are open and
retrocompact in $X$; it suffices to see that the complement of every
set of $\mathfrak{G}$ belongs to $\mathfrak{G}$. So let
$Z=\bigcup_{i\in I}(U_i\cap\complement{V_i})$, where $I$ is finite and $U_i$
and $V_i$ are open and retrocompact in $X$; we have
$\complement{Z}=\bigcap_{i\in I}(V_i\cup\complement{U_i})$, so
$\complement{Z}$ is a finite union of sets which are intersections of
some of the $V_i$ and some of the $\complement{U_i}$, hence
of the form
\oldpage[0\textsubscript{III}]{13}
$V\cap\complement{U}$, where $U$ is a union of some of the $U_i$ and
$V$ an intersection of some of the $V_i$; but we remarked
above that finite unions and intersections of open subsets retrocompact
in $X$ are open subsets retrocompact in $X$, whence the conclusion.

\begin{corollary}[9.1.4]
\label{0.9.1.4}
Every constructible subset of $X$ is retrocompact in $X$.
\end{corollary}

It suffices to show that if $U$ and $V$ are open and retrocompact in $X$,
then $U\cap\complement{V}$ is retrocompact in $X$; now, if $W$ is open and
quasi-compact in $X$, $W\cap U\cap\complement{V}$ is closed in the
quasi-compact space $W\cap U$, hence is quasi-compact.

In particular:

\begin{corollary}[9.1.5]
\label{0.9.1.5}
For an open subset $U$ of $X$ to be constructible, it is necessary and
sufficient that it be retrocompact in $X$. For a closed subset $F$ of
$X$ to be constructible, it is necessary and sufficient that the open subset
$\complement{F}$ be retrocompact.
\end{corollary}

\begin{env}[9.1.6]
\label{0.9.1.6}
An important case is that in which every quasi-compact open subset of $X$
is retrocompact, in other words, in which the intersection of two
quasi-compact open subsets of $X$ is quasi-compact
(cf. \hyperref[I.5.5.6]{I, 5.5.6}). When $X$ itself is
quasi-compact, this means that the open subsets retrocompact in
$X$ are identical with the quasi-compact open subsets of $X$, and the
constructible subsets of $X$ with the finite unions of sets of the form
$U\cap\complement{V}$, where $U$ and $V$ are open and quasi-compact.
\end{env}

\begin{corollary}[9.1.7]
\label{0.9.1.7}
For a subset of a Noetherian space $X$ to be constructible, it is necessary
and sufficient that it be a finite union of locally closed subsets of $X$.
\end{corollary}

\begin{proposition}[9.1.8]
\label{0.9.1.8}
Let $X$ be a topological space and $U$ an open subset of $X$.
\begin{enumerate}
  \item[(i)] If $T$ is a constructible subset of $X$, $T\cap U$ is a
    constructible subset of $U$.
  \item[(ii)] Suppose moreover that $U$ is retrocompact in $X$. For a
    subset $Z$ of $U$ to be constructible in $X$, it is necessary and
    sufficient that it be constructible in $U$.
\end{enumerate}
\end{proposition}

\begin{enumerate}
  \item[(i)] Using (\hyperref[0.9.1.3]{9.1.3}), we are reduced to
    showing that if $T$ is open and retrocompact in $X$, then $T\cap U$ is open
    and retrocompact in $U$, in other words, that for every quasi-compact open
    subset $W\subset U$, $T\cap U\cap W=T\cap W$ is quasi-compact, which
    follows immediately from the hypothesis.
  \item[(ii)] The condition being necessary by virtue of (i), it remains
    to prove that it is sufficient. Taking into account
    (\hyperref[0.9.1.3]{9.1.3}), it suffices to consider the case where $Z$
    is open and retrocompact \emph{in $U$}, since it will then follow that
    $U\setminus Z$ is constructible in $X$, and if $Z$, $Z'$ are two
    open subsets retrocompact \emph{in $U$},
    $Z\cap(U\setminus Z')$ will indeed be constructible in $X$. Now, if $W$ is
    open and quasi-compact in $X$ and $Z$ is open and retrocompact in $U$, we
    have $Z\cap W=Z\cap(W\cap U)$, and by hypothesis $W\cap U$ is open and
    quasi-compact in $U$; hence $W\cap Z$ is indeed quasi-compact, and
    consequently $Z$ is open and retrocompact in $X$, and \emph{a fortiori}
    constructible in $X$.
\end{enumerate}

\begin{corollary}[9.1.9]
\label{0.9.1.9}
Let $X$ be a topological space and $(U_i)_{i\in I}$ a finite covering of
$X$ consisting of open sets retrocompact in $X$. For a subset $Z$
of $X$ to be constructible in $X$, it is necessary and sufficient that for
every $i\in I$, $Z\cap U_i$ be constructible in $U_i$.
\end{corollary}

\begin{env}[9.1.10]
\label{0.9.1.10}
Suppose in particular that $X$ is quasi-compact and that every point
\oldpage[0\textsubscript{III}]{14}
of $X$ has a fundamental system of open neighbourhoods retrocompact
in $X$ (and \emph{a fortiori} quasi-compact); then the condition for a
subset $Z$ of $X$ to be constructible in $X$ is of a \emph{local} nature,
in other words, it is necessary and sufficient that for every $x\in X$
there exist an open neighbourhood $V$ of $x$ such that $V\cap Z$ is
constructible in $V$. Indeed, if this condition is satisfied, there
exists for every $x\in X$ an open neighbourhood $V$ of $x$ \emph{retrocompact
in $X$} and such that $V\cap Z$ is constructible in $V$, by virtue of the
hypothesis on $X$ and of (\hyperref[0.9.1.8]{9.1.8, (i)}); it then
suffices to cover $X$ by finitely many of these neighbourhoods and to apply
(\hyperref[0.9.1.9]{9.1.9}).
\end{env}

\begin{definition}[9.1.11]
\label{0.9.1.11}
Let $X$ be a topological space. We say that a subset $T$ of $X$ is
\emph{locally constructible in $X$} if, for every $x\in X$, there exists
an open neighbourhood $V$ of $x$ such that $T\cap V$ is constructible in $V$.
\end{definition}

It follows immediately from (\hyperref[0.9.1.8]{9.1.8, (i)}) that if $V$
is such that $V\cap T$ is constructible in $V$, then for every open subset
$W\subset V$, $W\cap T$ is constructible in $W$. If $T$ is locally
constructible in $X$, then, for every open subset $U$ of $X$, $T\cap U$ is
locally constructible in $U$, as follows from the preceding
remark. This same remark shows that the set of subsets locally
constructible in $X$ is stable under finite union and finite intersection;
it is clear, on the other hand, that it is also stable under passage to
complements.

\begin{proposition}[9.1.12]
\label{0.9.1.12}
Let $X$ be a topological space. Every set constructible in $X$ is
locally constructible in $X$. The converse is true if $X$ is
quasi-compact and its topology has a basis consisting of sets
retrocompact in $X$.
\end{proposition}

The first assertion follows from the definition
(\hyperref[0.9.1.11]{9.1.11}), and the second from
(\hyperref[0.9.1.10]{9.1.10}).

\begin{corollary}[9.1.13]
\label{0.9.1.13}
Let $X$ be a topological space whose topology has a basis consisting
of sets retrocompact in $X$. Then every subset $T$ locally
constructible in $X$ is retrocompact in $X$.
\end{corollary}

Indeed, let $U$ be a quasi-compact open set in $X$; $T\cap U$ is
locally constructible in $U$, hence constructible in $U$ by virtue of
(\hyperref[0.9.1.12]{9.1.12}), and consequently quasi-compact by virtue of
(\hyperref[0.9.1.4]{9.1.4}).

\subsection{Constructible sets in Noetherian spaces}
\label{subsection:0.9.2}

\begin{env}[9.2.1]
\label{0.9.2.1}
We have seen (\hyperref[0.9.1.7]{9.1.7}) that in a Noetherian space $X$, the
subsets constructible in $X$ are the \emph{finite unions of locally
closed subsets of $X$}.

The inverse image of a set constructible in $X$ under a continuous
map from a Noetherian space $X'$ to $X$ is constructible in $X'$.
If $Y$ is a constructible subset of a Noetherian space $X$, the subsets
of $Y$ which are constructible as subspaces of $Y$ are
identical with those which are constructible as subspaces of $X$.
\end{env}

\begin{proposition}[9.2.2]
\label{0.9.2.2}
Let $X$ be an irreducible Noetherian space and $E$ a constructible subset of
$X$. For $E$ to be everywhere dense in $X$, it is necessary and sufficient
that $E$ contain a non-empty open subset of $X$.
\end{proposition}

\oldpage[0\textsubscript{III}]{15}
The condition is evidently sufficient, every non-empty open set being
dense in $X$. Conversely, let
$E=\bigcup_{i=1}^n(U_i\cap F_i)$ be a constructible subset of $X$, the $U_i$
being non-empty open sets and the $F_i$ closed sets in $X$; we therefore have
$\overline{E}\subset\bigcup_i F_i$. Consequently, if $\overline{E}=X$, $X$ is
equal to one of the $F_i$, hence $E\supset U_i$, which completes the proof.

When $X$ has a generic point $x$
(\hyperref[0.2.1.2]{$0_{\mathrm I}$, 2.1.2}), the condition of
(\hyperref[0.9.2.2]{9.2.2}) is equivalent to the relation $x\in E$.

\begin{proposition}[9.2.3]
\label{0.9.2.3}
Let $X$ be a Noetherian space. For a subset $E$ of $X$ to be
constructible, it is necessary and sufficient that, for every irreducible
closed subset $Y$ of $X$, $E\cap Y$ be rare in $Y$ or contain a
non-empty open subset of $Y$.
\end{proposition}

The necessity of the condition comes from the fact that $E\cap Y$ must be a
constructible subset of $Y$, and from (\hyperref[0.9.2.2]{9.2.2}), since a
non-dense subset of $Y$ is necessarily rare in the irreducible space
$Y$ (\hyperref[0.2.1.1]{$0_{\mathrm I}$, 2.1.1}). To prove that the
condition is sufficient, let us apply the principle of Noetherian induction
(\hyperref[0.2.2.2]{$0_{\mathrm I}$, 2.2.2}) to the set
$\mathfrak{F}$ of closed subsets $Y$ of $X$ such that $Y\cap E$ is
constructible (relative to $Y$ or relative to $X$, which amounts to the
same thing): we may therefore suppose that for every closed subset $Y\ne X$
of $X$, $E\cap Y$ is constructible. Suppose first that $X$ is not
irreducible, and let $X_i$ ($1\leq i\leq m$) be its irreducible
components, which are necessarily finite in number
(\hyperref[0.2.2.5]{$0_{\mathrm I}$, 2.2.5}); by hypothesis, the
$E\cap X_i$ are constructible, hence so is their union $E$. Suppose
next that $X$ is irreducible; then, by hypothesis, either $E$ is
rare, so that $\overline{E}\ne X$ and $E=E\cap\overline{E}$ is constructible;
or $E$ contains a non-empty open subset $U$ of $X$, and is therefore the union
of $U$ and $E\cap(X\setminus U)$; but $X\setminus U$ is a closed set
distinct from $X$, so $E\cap(X\setminus U)$ is constructible; $E$
itself is consequently constructible, which completes the proof.

\begin{corollary}[9.2.4]
\label{0.9.2.4}
Let $X$ be a Noetherian space and $(E_\alpha)$ an increasing filtered
family of constructible subsets of $X$, such that:
\begin{enumerate}
  \item[$1^\circ$] $X$ is the union of the family $(E_\alpha)$.
  \item[$2^\circ$] Every irreducible closed subset of $X$ is contained in
    the closure of one of the $E_\alpha$.
\end{enumerate}
Then there exists an index $\alpha$ such that $X=E_\alpha$.

When every irreducible closed subset of $X$ has a generic point,
hypothesis $2^\circ$ can be omitted.
\end{corollary}

Let us apply the principle of Noetherian induction
(\hyperref[0.2.2.2]{$0_{\mathrm I}$, 2.2.2}) to the set
$\mathfrak{M}$ of closed subsets of $X$ contained in at least one of the
$E_\alpha$; we may therefore suppose that every closed subset $Y\ne X$ of $X$ is
contained in one of the $E_\alpha$. The proposition is evident if $X$ is not
irreducible, since each of the irreducible components $X_i$ of $X$
($1\leq i\leq m$) is contained in some $E_{\alpha_i}$, and there exists an
$E_\alpha$ containing all the $E_{\alpha_i}$. So suppose $X$
irreducible. By hypothesis, there exists $\beta$ such that
$X=\overline{E_\beta}$, hence (\hyperref[0.9.2.2]{9.2.2}) $E_\beta$
contains a non-empty open subset $U$ of $X$. But then the closed set
$X\setminus U$ is contained in some $E_\gamma$, and it suffices to take
$E_\alpha$ containing $E_\beta$ and $E_\gamma$. When every irreducible
closed subset $Y$ of $X$
\oldpage[0\textsubscript{III}]{16}
has a generic point $y$, there exists $\alpha$ such that $y\in E_\alpha$,
hence $Y=\overline{\{y\}}\subset\overline{E_\alpha}$, and condition
$2^\circ$ is a consequence of $1^\circ$.

\begin{proposition}[9.2.5]
\label{0.9.2.5}
Let $X$ be a Noetherian space, $x$ a point of $X$, and $E$ a constructible
subset of $X$. For $E$ to be a neighbourhood of $x$, it is necessary and
sufficient that for every irreducible closed subset $Y$ of $X$ containing $x$,
$E\cap Y$ be dense in $Y$ (if there exists a generic point $y$ of $Y$,
this also means (\hyperref[0.9.2.2]{9.2.2}) that $y\in E$).
\end{proposition}

The condition is evidently necessary; let us prove that it is sufficient.
Applying the principle of Noetherian induction to the set $\mathfrak{M}$
of closed subsets $Y$ of $X$ containing $x$ and such that $E\cap Y$ is a
neighbourhood of $x$ in $Y$, we may suppose that every closed subset $Y\ne X$
of $X$ containing $x$ belongs to $\mathfrak{M}$. If $X$ is not
irreducible, each of the irreducible components $X_i$ of $X$ containing
$x$ is distinct from $X$, so $E\cap X_i$ is a neighbourhood of $x$
relative to $X_i$; consequently, $E$ is a neighbourhood of $x$ in the union of
the irreducible components of $X$ containing $x$, and since this union is a
neighbourhood of $x$ in $X$, the same is true of $E$. If $X$ is
irreducible, $E$ is dense in $X$ by hypothesis, and therefore contains a
non-empty open subset $U$ of $X$ (\hyperref[0.9.2.2]{9.2.2}); the
proposition is then evident if $x\in U$; otherwise, $x$ is by hypothesis
interior to $E\cap(X\setminus U)$ relative to $X\setminus U$, so the closure in
$X$ of $X\setminus E$ does not contain $x$, and the complement of this
closure is a neighbourhood of $x$ contained in $E$, which completes the
proof.

\begin{corollary}[9.2.6]
\label{0.9.2.6}
Let $X$ be a Noetherian space and $E$ a subset of $X$. For $E$ to be an
open set in $X$, it is necessary and sufficient that for every irreducible
closed subset $Y$ of $X$ meeting $E$, $E\cap Y$ contain a non-empty
open subset of $Y$.
\end{corollary}

The condition is evidently necessary; conversely, if it is satisfied,
it implies that $E$ is constructible by virtue of
(\hyperref[0.9.2.3]{9.2.3}). Moreover,
(\hyperref[0.9.2.5]{9.2.5}) shows that $E$ is then a neighbourhood of each
of its points, whence the conclusion.

\subsection{Constructible functions}
\label{subsection:0.9.3}

\begin{definition}[9.3.1]
\label{0.9.3.1}
Let $h$ be a map from a topological space $X$ to a set $T$.
We say that $h$ is \emph{constructible} if $h^{-1}(t)$ is constructible for
every $t\in T$, and empty except for finitely many values of $t$; for
every subset $S$ of $T$, $h^{-1}(S)$ is then constructible. We say that $h$
is \emph{locally constructible} if every $x\in X$ has an open
neighbourhood $V$ such that $h|V$ is constructible.
\end{definition}

Every constructible function is locally constructible; the converse
is true when $X$ is quasi-compact and has a basis consisting of
open sets retrocompact in $X$ (in particular when $X$ is Noetherian).

\begin{proposition}[9.3.2]
\label{0.9.3.2}
Let $h$ be a map from a Noetherian space $X$ to a set $T$.
For $h$ to be constructible, it is necessary and sufficient that for every
irreducible closed subset $Y$ of $X$ there exist a non-empty subset $U$ of
$Y$, open relative to $Y$, on which $h$ is constant.
\end{proposition}

The condition is necessary: indeed, by hypothesis, $h$ takes on $Y$
only finitely many values $t_i$, and each of the sets
$h^{-1}(t_i)\cap Y$ is constructible in $Y$
(\hyperref[0.9.2.1]{9.2.1}); since they cannot all be rare subsets
of the space $Y$, at least one of them contains a non-empty open set
(\hyperref[0.9.2.3]{9.2.3}).

\oldpage[0\textsubscript{III}]{17}
To see that the condition is sufficient, let us apply the principle of
Noetherian induction to the set $\mathfrak{M}$ of closed subsets $Y$
of $X$ such that the restriction $h|Y$ is constructible; we may therefore
suppose that for every closed subset $Y\ne X$ of $X$, $h|Y$ is
constructible. If $X$ is not irreducible, the restriction of $h$ to
each of the irreducible components $X_i$ of $X$ (finite in number) is therefore
constructible, and it then follows immediately from the definition
(\hyperref[0.9.3.1]{9.3.1}) that $h$ is constructible. If $X$ is
irreducible, there exists by hypothesis a non-empty open subset $U$ of $X$
on which $h$ is constant; on the other hand, the restriction of $h$ to
$X\setminus U$ is constructible by hypothesis, and it follows immediately
that $h$ is constructible.

\begin{corollary}[9.3.3]
\label{0.9.3.3}
Let $X$ be a Noetherian space in which every irreducible closed subset
has a generic point. If $h$ is a map from $X$ to a set
$T$ such that, for every $t\in T$, $h^{-1}(t)$ is constructible, then $h$
is constructible.
\end{corollary}

Indeed, if $Y$ is a irreducible closed subset of $X$ and $y$ its generic
point, $Y\cap h^{-1}(h(y))$ is constructible and contains $y$, so
(\hyperref[0.9.2.2]{9.2.2}) this set contains a non-empty open subset of
$Y$, and it suffices to apply (\hyperref[0.9.3.2]{9.3.2}).

\begin{proposition}[9.3.4]
\label{0.9.3.4}
Let $X$ be a Noetherian space in which every irreducible closed subset
has a generic point, and $h$ a constructible map from $X$ to an
ordered set. For $h$ to be upper semi-continuous in $X$,
it is necessary and sufficient that for every $x\in X$ and every generization
(\hyperref[0.2.1.2]{$0_{\mathrm I}$, 2.1.2}) $x'$ of $x$, we have
$h(x')\leq h(x)$.
\end{proposition}

The function $h$ takes only finitely many values; saying that it is
upper semi-continuous therefore means that for every $x\in X$,
the set $E$ of $y\in X$ such that $h(y)\leq h(x)$ is a neighbourhood of $x$.
By hypothesis, $E$ is a constructible subset of $X$; on the other hand, saying
that a irreducible closed subset $Y$ of $X$ contains $x$ means that its
generic point $y$ is a generization of $x$; the conclusion then follows
from (\hyperref[0.9.2.5]{9.2.5}).

